\chapter{Tablice, obiekty, dane z JSON i model rekordow}

\section{Dlaczego dane zlozone sa potrzebne}
Prosta strona moze opierac sie na kilku pojedynczych zmiennych. Jednak prawdziwe projekty
pracuja zwykle na bogatszych strukturach: listach wpisow, produktach, autorach, kartach,
menu czy konfiguracjach. Do takich rzeczy MoonChunk daje dwa podstawowe narzedzia:

\begin{itemize}
  \item \mc{array} dla uporzadkowanych kolekcji,
  \item \mc{object} oraz \mc{dict} dla danych nazwanych.
\end{itemize}

\section{Tablice}
Tablica przechowuje elementy pod indeksami. Przyklad:

\begin{lstlisting}[style=moonchunk,caption={Prosta tablica}]
const colors: array = ["red", "green", "blue"];
print(colors[0]);
\end{lstlisting}

To naturalny wybor wtedy, gdy porzadek elementow ma znaczenie albo gdy chcesz iterowac
po danych petla z licznikiem.

\section{Czy tablica musi byc jednorodna}
MoonChunk jest tutaj dosc elastyczny. Tablica moze zawierac wartosci roznych typow,
choc w praktyce dla czytelnosci zwykle lepiej trzymac w niej elementy podobnego rodzaju.

\begin{lstlisting}[style=moonchunk,caption={Tablica mieszana}]
const mixed: array = [1, "ala", true];
\end{lstlisting}

Technicznie jest to mozliwe, ale projektowo warto zachowac umiar.

\section{Obiekty i slowniki}
Obiekt przechowuje pola po nazwie.

\begin{lstlisting}[style=moonchunk,caption={Prosty obiekt}]
const author: object = {
  name: "Alicja",
  role: "writer"
};
\end{lstlisting}

Slownik typu \mc{dict} dziala podobnie i jest szczegolnie wygodny, gdy chcesz wyraznie
zaakcentowac, ze pracujesz na strukturze danych klucz--wartosc.

\begin{lstlisting}[style=moonchunk,caption={Slownik dict}]
let point: dict = { x: 1, y: 2 };
\end{lstlisting}

\section{Dostep do pol przez kropke}
Najwygodniejszy i najczytelniejszy sposob pracy z obiektem to skladnia kropkowa.

\begin{lstlisting}[style=moonchunk,caption={Dostep do pola obiektu}]
print(author.name);
print(point.x);
\end{lstlisting}

To dobry domyslny styl, gdy nazwa pola jest znana na etapie pisania kodu.

\section{Dostep przez nawiasy kwadratowe}
Czasem klucz jest dynamiczny, np. przechowywany w innej zmiennej. Wtedy przydaje sie
skladnia indeksowana.

\begin{lstlisting}[style=moonchunk,caption={Dynamiczny dostep do pola}]
let key: string = "name";
print(author[key]);
\end{lstlisting}

Taki zapis jest bardzo mocny, bo pozwala programowi decydowac, po ktore pole ma siegnac.

\section{Klucze numeryczne i nietypowe}
MoonChunk pozwala tworzyc obiekty i slowniki takze z kluczami liczbowymi albo tekstowymi,
ktore nie sa zwyklymi identyfikatorami.

\begin{lstlisting}[style=moonchunk,caption={Nietypowe klucze obiektu}]
const details = {
  stats: { total: 3, labels: ["alpha", "bravo", "charlie"] },
  "meta-key": "enabled",
  101: "numeric-key"
};
\end{lstlisting}

To przydaje sie przy pracy z danymi przychodzacymi z zewnatrz, gdzie klucze nie zawsze
sa idealnie dopasowane do skladni zwyklych identyfikatorow.

\section{Modyfikowanie pol i elementow}
MoonChunk pozwala nie tylko odczytywac dane z obiektow i tablic, ale tez je zmieniac.

\begin{lstlisting}[style=moonchunk,caption={Modyfikacja danych zlozonych}]
details.stats.total = details.stats.total + 1;
details.stats.labels[1] = "beta";
\end{lstlisting}

To bardzo wygodne, gdy chcesz przetworzyc dane przed wygenerowaniem finalnego HTML-a.

\section{Tablice i obiekty jako model rekordow}
MoonChunk nie wprowadza osobnej skladni typu \emph{krotka} albo \emph{rekord}. W praktyce
rolę rekordu z nazwanymi polami pelni obiekt albo \mc{dict}.

Przyklad punktu jako rekordu:

\begin{lstlisting}[style=moonchunk,caption={Rekord modelowany przez obiekt}]
const point: object = { x: 1, y: 2 };
\end{lstlisting}

Jesli potrzebujesz nazwanych pol, to zwykle jest wlasnie najlepsze rozwiazanie.
Jesli potrzebujesz uporzadkowanej listy bez nazw pol, uzyj tablicy.

\section{Wczytywanie danych z pliku JSON przez \mc{data(...)}}
To jedna z najpraktyczniejszych funkcji w MoonChunk. Pozwala pobrac dane z pliku JSON
i pracowac na nich jak na zwyklych strukturach danych.

\begin{lstlisting}[style=moonchunk,caption={Wczytanie JSON do zmiennej}]
const posts = data("./data/posts.json");
\end{lstlisting}

Po takim zapisie \mc{posts} staje sie wartoscia, z ktora mozna pracowac w petlach,
warunkach i interpolacjach.

\section{Przykladowy plik JSON}
\begin{lstlisting}[style=basecode,caption={Przykladowe dane posts.json}]
[
  {
    "title": "Pierwszy wpis",
    "slug": "pierwszy-wpis",
    "published": true,
    "views": 120
  },
  {
    "title": "Drugi wpis",
    "slug": "drugi-wpis",
    "published": false,
    "views": 30
  }
]
\end{lstlisting}

Taki plik mozesz potem przeglądać w MoonChunk przez indeks i pola obiektu.

\section{Przyklad pracy z danymi JSON}
\begin{lstlisting}[style=moonchunk,caption={Filtrowanie i wyswietlanie wpisow}]
const posts = data("./data/posts.json");

page "" {
  for (let i: int = 0; i < 2; i++) {
    if (posts[i].published) {
      content {
        <article>
          <h2>{posts[i].title}</h2>
          <p>Views: {posts[i].views}</p>
        </article>
      };
    };
  };
};
\end{lstlisting}

To prosty, ale bardzo realny scenariusz: część danych jest pokazywana tylko wtedy,
gdy dany wpis jest opublikowany.

\section{Kiedy uzywac \mc{object}, a kiedy \mc{dict}}
Z perspektywy czystego uzytkowania oba typy sa do siebie bardzo podobne. Warto przyjac
prostą konwencje:

\begin{itemize}
  \item \mc{object} -- gdy myslisz o konkretnej strukturze modelu,
  \item \mc{dict} -- gdy myslisz o slowniku danych i kluczach.
\end{itemize}

Nie jest to sztywna roznica filozoficzna, ale bardzo pomaga w czytelnosci projektu.

\section{Typowe bledy przy danych zlozonych}
\begin{itemize}
  \item probowanie odczytu pola z \mc{null} albo \mc{undefined},
  \item niepoprawny typ indeksu w nawiasach kwadratowych,
  \item zakladanie, ze dana struktura ma pole, ktorego nie ma,
  \item probowanie traktowania prostej liczby jak obiektu,
  \item wpisywanie skomplikowanej logiki danych bezposrednio do HTML-a zamiast przygotowac wartosci wczesniej.
\end{itemize}

\section{Dobra praktyka: dane najpierw, HTML potem}
Przy pracy z tablicami i obiektami bardzo oplaca sie najpierw przygotowac czyste dane,
a dopiero potem budowac z nich tresc strony. To redukuje balagan i sprawia, ze \mc{content}
pozostaje glownie warstwa prezentacji, a nie mieszanina logiki i znacznikow.
